\clearpage\fancyhead[L]{\small PALLAS.}\SourceLeaf{381}
\begin{center}\textbf{PALLAS.}\par\vspace{4pt}\hrule\vspace{4pt}\fontsize{10}{12}\selectfont Astronomisches Jahrbuch für 1817. Seite 212\;..\;215. Berlin 1814.\par\vspace{4pt}\hrule\vspace{4pt}\end{center}

\begin{flushright}
Göttingen, den 29. Mai 1814.
\end{flushright}

\par\indent Mein verehrtester Freund, ich schulde Ihnen noch meinen herzlichsten Dank für das wertvolle Geschenk Ihres astronomischen Jahrbuchs für 1816. Ich habe meinen Dank aufgeschoben, weil ich ihn mit einigen Beiträgen für das neue Jahrbuch verbinden wollte. Die Arbeit an dem, was ich Ihnen über die letzte Vestaopposition schicken wollte, hat sich bisher jedoch verzögert. Ich werde Ihnen das aber noch nachreichen.

\Needspace{14\baselineskip}\par\indent Über die letzte Junoopposition habe ich Ihnen zwar schon einiges mitgeteilt. Da ich jetzt noch Verschiedenes hinzufügen kann und nicht mehr genau weiß, wie viel ich Ihnen bereits geschrieben habe, will ich hier alles Wesentliche zusammenstellen. Das Wetter erlaubte mir um die Zeit der Opposition nur eine Beobachtung, die ich allerdings für sehr gut halte. Erst fast drei Wochen später konnte ich zur Kontrolle eine zweite machen, die ich nach der Reduktion von Herrn \textsc{Nicolai} beifüge:

\begin{center}\fontsize{10}{12}\selectfont\setlength{\tabcolsep}{3pt}\renewcommand{\arraystretch}{1.15}\begin{tabular}{l|c|rrr|rrr}
1813 & Mittlere Zeit & \multicolumn{3}{c|}{Scheinbare Rektasz.} & \multicolumn{3}{c}{Südl. Dekl.}\\\hline
Nov. 19 & $9^{\mathrm u}\;46^{\mathrm m}\;46^{\mathrm s}$ & $60^\circ$ & $30'$ & $35''0$ & $3^\circ$ & $5'$ & $9''3$\\
Dez. 9 & $7\;9\;38$ & 56 & 56 & 47.4 & 3 & 56 & 35.3\\
\end{tabular}\end{center}

\Needspace{14\baselineskip}\par\indent Ich hatte die vorläufige Berechnung der Opposition zwar bereits auf meine erste Beobachtung gestützt. Da ich aber zum einen nur flüchtig gerechnet hatte und zum anderen \textsc{Bessel} mir danach noch seine Meridianbeobachtungen mitteilte (die Sie zweifellos von ihm selbst erhalten haben werden), bat ich einen geschickten und talentierten Hörer meiner Vorlesungen, Herrn \textsc{Möbius}, die Ergebnisse entsprechend zu berichtigen. Er findet für die Opposition:

\begin{center}19. November 1813. $18^{\mathrm u}\;12^{\mathrm m}\;0^{\mathrm s}$ Mittlere Zeit in Göttingen\end{center}
\begin{center}\fontsize{10}{12}\selectfont\setlength{\tabcolsep}{2pt}\renewcommand{\arraystretch}{1.15}\begin{tabular}{lcr}
Wahre Länge & $.\;.\;.\;.\;.\;.\;.\;.\;.\;.\;.\;.\;.\;.$ & $57^\circ\;34'\;2''1$\\
Geozentrische Breite & $.\;.\;.\;.\;.\;.\;.\;.\;.\;.\;.$ & $23\;18\;46.1\;\mathrm{S.}$\\
\end{tabular}\end{center}

\Needspace{14\baselineskip}\par\indent Zugleich hat er die Elemente an diese und die drei vorhergehenden Oppositionen von 1810, 1811 und 1812 angepasst und folgendes Ergebnis erhalten:

\begin{center}\fontsize{9}{11}\selectfont\setlength{\tabcolsep}{2pt}\renewcommand{\arraystretch}{1.15}\begin{tabular}{lcr}
Epoche 1810, Meridian von Göttingen & $.\;.\;.\;.\;.\;.\;.$ & $95^\circ\;29'\;53''2$\\
Tägliche mittlere tropische Bewegung & $.\;.\;.\;.\;.\;.\;.$ & $812''7140$\\
Länge des Perihels (1810 sid. ruh.) & $.\;.\;.\;.\;.\;.\;.$ & $53^\circ\;6'\;43''0$\\
Länge des aufsteigenden Knotens (1810) & $.\;.\;.\;.\;.\;.$ & $171\;6\;45.0$\\
Exzentrizitätswinkel & $.\;.\;.\;.\;.\;.\;.\;.\;.\;.\;.\;.\;.$ & $14\;43\;9.5$\\
Neigung der Bahn & $.\;.\;.\;.\;.\;.\;.\;.\;.\;.\;.\;.\;.\;.$ & $13\;4\;12.9$\\
Logarithmus der großen Halbachse & $.\;.\;.\;.\;.\;.\;.$ & $0.4267616$\\
\end{tabular}\end{center}

\par\indent Außerdem hat dieser geschickte Rechner für die nächste Erscheinung der Juno eine Ephemeride berechnet, von der ich Ihnen gern eine Abschrift beilege.

\SourceLeaf{382}\par\indent Für Pallas hat Herr \textsc{Nicolai} (der seit einem halben Jahr als Gehilfe an der Seeberger Sternwarte angestellt ist) ebenfalls eine Ephemeride berechnet und mir die beiliegende Kopie für Ihr Jahrbuch geschickt. Zum Auffinden des Planeten (der diesmal ohnehin ziemlich hell sein wird) ist sie gewiss genauer als nötig, obwohl sie nach den früheren Elementen berechnet wurde und deren Veränderung durch die fortwirkenden Störungen des Jupiters unberücksichtigt blieb. Es ist aber doch interessant, den Einfluss dieser Störungen \emph{im Voraus} anzugeben. Deshalb habe ich Herrn \textsc{Möbius} diese ziemlich mühsame Arbeit übertragen. Das Ergebnis ist, dass die Rektaszension in der Ephemeride von Herrn \textsc{Nicolai} um die Zeit der Opposition um 1$\frac{1}{2}$ Minuten vermindert werden muss, während sich die Deklination nur um wenige Sekunden vermindert. Für die Opposition selbst erhält er im Voraus Folgendes:

\begin{center}25. Oktober 1814. $12^{\mathrm u}\;39^{\mathrm m}\;50^{\mathrm s}$ Göttinger Zeit.\end{center}
\begin{center}\fontsize{10}{12}\selectfont\setlength{\tabcolsep}{2pt}\renewcommand{\arraystretch}{1.15}\begin{tabular}{lcr}
Wahre Länge & $.\;.\;.\;.\;.\;.\;.\;.\;.\;.$ & $31^\circ\;58'\;28''4$\\
Heliozentrische Breite & $.\;.\;.\;.\;.\;.\;.$ & $23\;38\;32.7\;\mathrm{S.}$\\
\end{tabular}\end{center}

\noindent (Die geozentrische Breite hat er mir mitzuteilen vergessen. Da er sich jetzt nicht mehr hier aufhält, kann ich sie zum Vergleich mit dem Ergebnis von Herrn \textsc{Nicolai} nicht beifügen.)

\par\indent Ich habe jetzt angefangen, auch die Störungen der Pallas durch Mars zu berechnen. Sie werden zwar fast alle sehr klein, aber doch zahlreich sein. Die Gleichungen für Jupiter, Saturn und Mars, die über $0''1$ liegen, werden insgesamt mehr als 1000 an der Zahl sein.

\Needspace{10\baselineskip}\par\indent Vor ein paar Tagen habe ich aus München ein ausgezeichnetes Heliometer mit 43 Zoll Brennweite und 34 Linien Öffnung erhalten. Der geniale Instrumentenbauer hat unter anderem die wichtige Verbesserung vorgenommen, dass er es für wiederholte Messungen eingerichtet hat, indem er beide Objektivhälften unabhängig voneinander beweglich gemacht hat. Dadurch vermeidet man den nachteiligen Einfluss kleiner örtlicher Unebenheiten der Schraube vollständig. Die Präzision des Fernrohrs ist, wie man es von \textsc{Frauenhofer}s Arbeiten gewohnt ist, außerordentlich hoch.
